\section*{Capítulo \ref{ch:g1:my-body} --- Mi cuerpo}

\begin{solution}{exo:g1:my-body:1}
Cabeza, tronco y \omterm{def:g1:my-body:parts}{extremidades} (dos \omterm{def:g1:my-body:parts}{brazos} y dos \omterm{def:g1:my-body:parts}{piernas}). El cuello
une la cabeza al tronco.
\end{solution}

\begin{solution}{exo:g1:my-body:2}
Un \omterm{def:g1:my-body:parts}{brazo} termina en una mano; una \omterm{def:g1:my-body:parts}{pierna} termina en un pie.
\end{solution}

\begin{solution}{exo:g1:my-body:3}
(a) El \omterm{def:g1:my-body:joint}{codo}. (b) La \omterm{def:g1:my-body:joint}{rodilla}. (c) La \omterm{def:g1:my-body:joint}{muñeca}.
\end{solution}

\begin{solution}{exo:g1:my-body:4}
El \omterm{def:g1:my-body:joint}{hombro} giró el \omterm{def:g1:my-body:parts}{brazo} hacia el otro lado, el \omterm{def:g1:my-body:joint}{codo} lo dobló y la
\omterm{def:g1:my-body:joint}{muñeca} inclinó la mano hasta la oreja.
\end{solution}

\begin{solution}{exo:g1:my-body:5}
Se parecen: el mismo plano general --- una cabeza, un cuerpo (tronco),
cuatro \omterm{def:g1:my-body:parts}{extremidades}. Se diferencian: tú te sostienes erguido sobre dos
\omterm{def:g1:my-body:parts}{piernas} con los dos \omterm{def:g1:my-body:parts}{brazos} libres, mientras que el gato anda a cuatro
patas; y el gato tiene cola, y tú no.
\end{solution}

\begin{solution}{exo:g1:my-body:6}
Mirando sin más, adivinas la coronilla de reojo y la marca queda
demasiado alta o demasiado baja. El libro plano, en horizontal y
tocando la pared, lleva lo más alto de la cabeza justo hasta la pared:
la misma medida justa todas las veces.
\end{solution}

\begin{solution}{exo:g1:my-body:7}
Enseñan que tu cuerpo se hizo más alto entre las dos fechas. No ---
crecer es demasiado lento para notarlo; solo las marcas lo hacen visible.
\end{solution}

\begin{solution}{exo:g1:my-body:8}
Respuesta abierta: andar bien, sentarse, coger cosas, aplaudir,
escribir --- casi todos los movimientos resultan usar \omterm{def:g1:my-body:joint}{rodillas},
caderas, \omterm{def:g1:my-body:joint}{codos}, \omterm{def:g1:my-body:joint}{hombros} o \omterm{def:g1:my-body:joint}{muñecas}. Sin articulaciones, el cuerpo se
movería como un palo tieso.
\end{solution}

\begin{solution}{exo:g1:my-body:9}
Que cada cuerpo es normal a su tamaño y a su ritmo: unos niños crecen
antes y otros después, y todos llegan a su propio tamaño de persona
mayor. Las marcas de altura son para verte crecer a ti, no para echar
una carrera con nadie.
\end{solution}

\begin{solution}{exo:g1:my-body:10}
Corta en las caderas (para que las \omterm{def:g1:my-body:parts}{piernas} se doblen hacia delante y
pueda sentarse), en las \omterm{def:g1:my-body:joint}{rodillas} (para que las pantorrillas cuelguen)
y se sienta aún mejor con un corte en el cuello y en los \omterm{def:g1:my-body:joint}{hombros} para
inclinarse y equilibrarse. Los cortes copian las articulaciones del
cuerpo: cadera, \omterm{def:g1:my-body:joint}{rodilla}, cuello, \omterm{def:g1:my-body:joint}{hombro}.
\end{solution}
